\section{The Beam 1--Beam 2 interface}

The smallest defensible interface is
\begin{equation}
(F_t,S_t)
\xrightarrow{\;\mathcal R\;}
a_t
\xrightarrow{\;\mathcal W\;}
w_t
\xrightarrow{\;\mathcal E\;}
e_t
\xrightarrow{\;\mathcal L\;}
S_{t+1},
\label{eq:interface}
\end{equation}
where \(S_t\) is an intentionally generic state containing whatever current knowledge and
competence variables are required.

The use of \(S_t\), rather than prematurely collapsing everything into one scalar competence,
is important. A future model may require at least two components:
\begin{equation}
S_t=(K_t,C_t),
\end{equation}
where \(K_t\) represents the availability of problem-solution knowledge and \(C_t\) represents
proficiency or efficiency. The current theoretical foundations do not yet justify eliminating
either component.

Here \(w_t\) is real operational work and \(e_t\) is learning-relevant exposure. They may
coincide in a source model or special HLS instantiation, but they are not identified by
default. The maps \(\mathcal W\), \(\mathcal E\), and \(\mathcal L\) must state their
types, units, domains, and assumptions.

\subsection{A task-indexed learning abstraction}

If attention is restricted to the intensive competence margin, a simple representation
consistent with the mechanisms of Beam~2 is
\begin{align}
E_{t+1}(q) &= \rho_q E_t(q) + e_t(q),\\
C_t(q) &= L_q(E_t(q)),
\label{eq:experience}
\end{align}
where \(E_t(q)\) is retained accumulated exposure, \(e_t(q)\) is current
learning-relevant exposure, \(\rho_q\) is a retention factor and \(L_q\) is a learning
function.

This is \textbf{not} presented as an established equation from Arrow, Argote, Gibbons--Waldman
or Gutjahr. It is a minimal synthesis that preserves their relevant mechanisms: experience,
task specificity, learning and depreciation.

\subsection{Source support and remaining mapping burden}

The interface in \eqref{eq:interface} requires only two substantive claims.

First, operational organization determines which tasks are performed or escalated. This is
the subject of Beam~1.

Second, task-performance experience can alter future knowledge or competence. This is
supported independently by Beam~2 and is made particularly explicit by the experience-to-
knowledge cycle of \citet{argotemiron2011} and the dynamic competence equation of
\citet{gutjahr2011}.

These source results make the scaffold plausible, but they do not establish the particular
maps or their semantic equivalence in HLS. The bridge remains \textbf{CANDIDATE /
HYPOTHESIS}; integration must proceed through the canonical ontology rather than direct
paper-to-paper variable identification.
